\section{Hukum Lanjutan: Pertukaran, Distributif, dan De Morgan}

Setelah melengkapi diri berbekal tameng dasar penangkal Redundansi dan Identitas di subbab lampau, petualangan aljabar menyuguhkan pisau operasi yang sanggup merobek-robek urutan kompilasi struktural kurung dan memecah belah operator \cite{rosen2019}. Inilah ranah unjuk kekuatan tri-triumvirat Komutatif, Distributif silang, serta palu godam De Morgan.

\subsection{1. Memutar Urutan dan Menyisipkan Kurung (Komutatif \& Asosiatif)}

Sebagaimana sifat kodrat hakiki penjumlahan $(3+2 = 2+3)$ dan pengelompokan perkalian pada aljabar angka riil konvensional, relasi dasar Logika Boolean (AND, OR) patuh membatu membebek pada tatanan sakral linier yang membebaskannya berotasi susunan ruang.

\begin{itemize}
    \item \textbf{Hukum Komutatif (Pertukaran Barter Bebas)} \\
    Menegaskan bahwa urutan sapaan parameter input sama sekali dibebaskan tanpa beban kasta diskriminasi. Huruf di kiri halal dihimpit ke arah seberang kanan.
    \begin{align*}
        p \vee q &\equiv q \vee p \\
        p \wedge q &\equiv q \wedge p 
    \end{align*}

    \item \textbf{Hukum Asosiatif (Pergeseran Skala Kurung)} \\
    Sepanjang deretan pipa argumen **mengalirkan operator pertalian kasta yang SEJENIS mutlak 100\% seragam murni** (baik deretan majemuk full OR, maupun formasi rakitan murni AND), batas partisi isolasi sekat tanda kurung dapat dibongkar pasang bergeser tempat.
    \begin{align*}
        (p \vee q) \vee r &\equiv p \vee (q \vee r) \\
        (p \wedge q) \wedge r &\equiv p \wedge (q \wedge r) 
    \end{align*}
\end{itemize}

\subsection{2. Hukum Silang Pedang: Distributif}

Ketika barisan selang pipa argumen menjumpai persimpangan sekat kurung di mana kasta operator di luar pagar \textbf{LALU SANGAT BERBEDA WUJUD SECARA KODRAT} (misalkan pihak \textit{Or} di luar memanjat pagar gerbang benteng kurung pihak \textit{And}), terjadilah invasi hukum pemecahan Distribusi Perkalian silang masuk ke urat nadi kedua penghuni kurung tujuannya.

Hukum Distributif Logika beroperasi cerminan seidentik paripurna serupa silang tebaran perkalian mendistribusi ke ruas tambahan aljabar $x \cdot (y+z) = x\cdot y + x\cdot z$.

\begin{align*}
    p \vee (q \wedge r) &\equiv (p \vee q) \wedge (p \vee r) \\
    p \wedge (q \vee r) &\equiv (p \wedge q) \vee (p \wedge r) 
\end{align*}
Mengingat tabiat hukumnya berlaku fleksibel dua kompas, kemampuan menilik dari ruas KANAN panjang terurai gila-gilaan, lalu secara jeli insting Anda memfaktorkannya mengerucut ringkas melebur menembus wujud kurung ruas KIRI adalah sejati parameter kualifikasi tertinggi kelulusan mahasiswa merengkuh master kelas Logika.

\subsection{3. Tangan Besi Pemecah Kurung: Hukum De Morgan}

Satu-satunya algojo matematika terkejam yang memperoleh izin yudisial membacok menginfiltrasi ke jantung kurung pertahanan lawan, meledakkan kepingan properti, lalu menyuntikkan bisa \textit{Reverse} mutasi memutarbalikkan kodrat operator AND dirubah radikal menjadi monster OR, dan sebaliknya. Dinobatkan atas jerih payah peletak batu fondasi Augustus De Morgan \cite{wiki_first_order_logic}.

\begin{teorema}[Hukum Pembalik De Morgan]
Masuknya rudal negasi $\neg$ menghunjam dari luar membran pagar kurung $(p \,??\, q)$ tak cukup merubah kawan variabel menitis lawan $\neg p$ maupun $\neg q$, serbuan rudal tersebut mutlak menyabet simbol kompas membalik vertikal orientasi operator tengah mereka! Pihak $\land$ berontak membias mendongak menitis $\lor$.
\end{teorema}

\begin{align*}
    \neg(p \wedge q) &\equiv \neg p \vee \neg q \\
    \neg(p \vee q) &\equiv \neg p \wedge \neg q
\end{align*}

Menguasai teknik merangkai baris manipulasi Distributif dan injeksi tembak balik De Morgan adalah jalan ninja *Cheat Mode* menyederhanakan rangkaian papan sirkuit IC di masa transisi mahasiswa teknik kelak.
